\begin{proof}[Author proof of Theorem9]
The proof follows several steps.
\paragraph{STEP1}
Since $f$ satisfies the left one-to-one condition, by Proposition6, $B$ computes a one-to-one mapping; hence, we can extend it to a permutation $\tau:\{-1,1\}^{n+k}\to\{-1,1\}^{n+k}$ as follows:
\[\tau(y_1,\ldots,y_{n+k})=\begin{cases}B(y_1,\ldots,y_n)&\text{if }y_{n+1}=\cdots=y_{n+k}=1,\\\text{arbitrary choice}&\text{otherwise},\end{cases}\]
where by arbitrary choice we mean one of the $(2^{n+k}-2^n)!$ possible ways of completing the definition so as to yield a permutation. Define $\sigma=\tau^{-1}$ and note that $\sigma$ is a permutation as well.
\paragraph{STEP2}
For every $x\in\{-1,1\}^n$ and $R\subseteq\{n+1,\ldots,n+k\}$, define $f_x^R:\{-1,1\}^{n+k}\to\{-1,1\}^{n+k}$ by
\[f_x^R(y_1,\ldots,y_{n+k})=\begin{cases}f_x(y_1,\ldots,y_n)&\text{if }y_{n+1}=\cdots=y_{n+k}=1,\\\chi_{S(x)}(y_1,\ldots,y_n)\cdot\chi_R(y_{n+1},\ldots,y_{n+k})&\text{otherwise}.\end{cases}\]
What can we say about these functions?
\begin{itemize}
\item Fix $x\in\{-1,1\}^n$, and denote by $C_x$ the circuit obtained from $C$ when Alice is given $x$ as input. Now, consider $y\in\{-1,1\}^{n+k}$.
\begin{itemize}
\item If $y_{n+1}=\cdots=y_{n+k}=1$, then $f_x^R(y)=f_x(y)$ by definition; hence, by the correctness of $P$ and the definition of $\sigma$, we have that $f_x^R$ agrees with $C_x\circ\sigma^{-1}$ on all such $y$'s.
\item Otherwise, let $i\in\{n+1,\ldots,n+k\}$ such that $y_i=-1$. For every $R\subseteq\{n+1,\ldots,n+k\}$ that contains $i$, we have that $C_x\circ\sigma^{-1}$ agrees with exactly one of $f_x^R,f_x^{R\setminus\{i\}}$ on the input $(y_1,\ldots,y_{n+i-1},-1,y_{n+i+1},\ldots,y_{n+k})$; thus, for exactly half the subsets $R\subseteq\{n+1,\ldots,n+k\}$, $f_x^R$ agrees with $C_x\circ\sigma^{-1}$ on $y$. Therefore,
\[\Pr_{R\sim2^{[n+k]\setminus[n]}}[f_x^R(y)=C_x(\sigma^{-1}(y))]=\frac12.\]
This holds for any $y$ such that $(y_{n+1},\ldots,y_{n+k})\ne(1,\ldots,1)$; hence,
\[\underset{\substack{y\sim\{-1,1\}^{n+k}\\\exists j\in[k]:y_{n+j}=-1}}{\mathbb E}\left[\Pr_{R\sim2^{[n+k]\setminus[n]}}[f_x^R(y)=C_x(\sigma^{-1}(y))]\right]=\frac12,\]
which implies there exists $R(x)\subseteq[n+k]\setminus[n]$ such that
\[\Pr_{\substack{y\sim\{-1,1\}^{n+k}\\\exists j\in[k]:y_{n+j}=-1}}[f_x^{R(x)}(y)=C_x(\sigma^{-1}(y))]\geq\frac12.\]
\end{itemize}
It follows that the fraction of inputs on which $f_x^{R(x)}$ and $C_x\circ\sigma^{-1}$ agree is at least
\[\frac{2^n+(1/2)\cdot(2^{n+k}-2^n)}{2^{n+k}}=\frac12+\frac1{2^{k+1}},\]
which is also the fraction of inputs on which $F_x^{R(x)}\triangleq f_x^{R(x)}\circ\sigma$ and $C_x$ agree.
\item The second thing we observe is that for any $x\ne x'\in X$,
\begin{align*}\langle F_x^{R(x)},F_{x'}^{R(x')}\rangle&=\langle f_x^{R(x)}\circ\sigma,f_{x'}^{R(x')}\circ\sigma\rangle=\langle f_x^{R(x)},f_{x'}^{R(x')}\rangle\\&=\mathbb E_{y\sim\{-1,1\}^{n+k}}[f_x^{R(x)}(y)\cdot f_{x'}^{R(x')}(y)]\\&=\mathbb E_{y\sim\{-1,1\}^n}[f_x(y)\cdot f_{x'}(y)]\cdot2^{-k}\\&\quad+\sum_{\substack{z\in\{-1,1\}^k\\\exists j\in[k]:z_j=-1}}\chi_{R(x)}(z)\cdot\chi_{R(x')}(z)\cdot\mathbb E_{y\sim\{-1,1\}^n}[\chi_{S(x)}(y)\cdot\chi_{S(x')}(y)]\cdot2^{-k}\\&=\langle f_x,f_{x'}\rangle\cdot2^{-k}\\&\quad+\sum_{\substack{z\in\{-1,1\}^k\\\exists j\in[k]:z_j=-1}}\chi_{R(x)}(z)\cdot\chi_{R(x')}(z)\cdot\underbrace{\langle\chi_{S(x)},\chi_{S(x')}\rangle}_{0}\cdot2^{-k}\\&=\langle f_x,f_{x'}\rangle\cdot2^{-k},\end{align*}
which implies
\begin{align*}\mathbb E_{x\ne x'\sim X}[\langle F_x^{R(x)},F_{x'}^{R(x')}\rangle^2]&\leq2^{-2k}\cdot\mathbb E_{x\ne x'\sim X}[\langle f_x,f_{x'}\rangle^2]\\&\underset{\text{assumption}}\leq2^{-2k}\cdot\frac{2^{2k}}{36|X|^2}=\frac1{36|X|^2}.\end{align*}
\end{itemize}
\paragraph{STEP3}
Let $V$ be the inner product space of all functions $\{-1,1\}^{n+k}\to\mathbb R$, and let $U\leq V$ be the subspace of all functions of degree up to $t$, which is spanned by $\{\chi_T\}_{|T|\leq t}$ and has dimension $\dim U=\binom{n+k}{\leq t}$. By the Projection Lemma, there exists $x^*\in X$ such that
\[\|(F_{x^*}^{R(x^*)})^{\leq t}\|^2\leq\frac{13\binom{n+k}{\leq t}}{|X|}\ \underset{\text{assumption}}\leq\ \frac1{2^{2(k+1)}}\quad\Longrightarrow\quad\|(F_{x^*}^{R(x^*)})^{\leq t}\|\leq\frac1{2^{k+1}}.\]
Since $\|(F_{x^*}^{R(x^*)})^{\leq t}\|\leq1/2^{k+1}$ and $C_{x^*}$ is an unbounded fan-in circuit of depth $h$ and size $\leq M$ that agrees with $F_{x^*}^{R(x^*)}$ on at least $1/2+1/2^{k+1}$ fraction of inputs, by Proposition18,
\[M\geq2^{\Omega_h([t/(1-2\log2^{-(k+1)})]^{1/(h-1)})}=2^{\Omega_h([t/(2k+3)]^{1/(h-1)})}=2^{\Omega_h([t/k]^{1/(h-1)})}.\]
\end{proof}
